% =====================================================================================================
% Kapitel 5: Umstellung des QR-Code-Scanners auf Ethernet                       Zielumfang: ca. 6 Seiten
% =====================================================================================================

\section{Umstellung des QR-Code-Scanners auf Ethernet}

\subsection{Ist-Analyse der bestehenden Anbindung}
% INHALT:  - Matrix 200 + CBX100, RS232 an Laptop; Scanner.py (pyserial, COM3, 19200, 8N1, „LON“)
%          - Schwachstellen: Teilcodes bei Timeout, kein stop_event, kein Reconnect, Ziffernfilter
% QUELLEN: Code T3200 (als Quelle angeben), Handbuch Matrix 200 / CBX100
% UMFANG:  ca. 1 Seite

\subsection{Variantenvergleich}
% INHALT:  - Tabelle: virtueller COM-Port vs. TCP-Socket (vs. ggf. Ethernet direkt am Scanner,
%            falls Matrix-200-Variante das kann) – Kriterien: Aufwand, Robustheit, Treiber, Lehrwert
%          - begründete Entscheidung
% UMFANG:  ca. 1 Seite

\subsection{Netzwerkkonzept}
% INHALT:  - Adressplan (Tabelle), Topologie-Abbildung (Switch, SPS, HMI, PC, Konverter, AL1100)
%          - Port, Betriebsart des Konverters
% UMFANG:  ca. 0,5 Seiten

\subsection{Konfiguration von Scanner und Konverter}
% INHALT:  - serielle Parameter (Tabelle), Terminator (gemessen!), Start-String
%          - Anschluss Konverter an CBX100 (Klemmen laut Handbuch)
% UMFANG:  ca. 1 Seite

\subsection{Anpassung der Python-Anwendung}
% INHALT:  - neue Scanner-Klasse: Socket, Puffer/Terminator, Timeout, Reconnect, stop_event
%          - Schnittstelle Scanner() unverändert → restliches Programm unberührt
%          - Listing (\pythonstyle) mit den wesentlichen Teilen, kein kompletter Code im Text
% UMFANG:  ca. 1,5 Seiten

\subsection{Test und Inbetriebnahme}
% INHALT:  - Testfälle als Tabelle: Einzelscan, Dauerbetrieb, Kabel ziehen/stecken, falscher Code,
%            Neustart Konverter, Gesamttest mit SPS
% UMFANG:  ca. 0,5 Seiten

\subsection{Fazit}
% INHALT:  Ergebnis bezogen auf Ziel 2, Überleitung zu Kapitel 6
